\chapter{共同余数障碍与锐下界}\label{sec:obstruction}

上一讲给出的 $G$ 看似只是对差分取最大公因子，本讲解释它为什么抓住了所有必要的算术条件。
先固定一个素数 $\ell\ne p$，看全部数模 $\ell$ 的余数；
再在阈值下方考察分母中 $p$ 的指数。前者证明“有障碍就不能成功”，
后者用来排除反例中可能存在的分数绕路。
\section{归一化为什么不改变可达性}

\begin{lemma}[归一化与操作序列相容]\label{lem:affine}
任一平均词 $W$ 满足 $W(a\mathbf1+bx)=a\mathbf1+bWx$。因此配置达到均值，等价于其本原中心化代表达到零。
\end{lemma}
\begin{proof}
每个 $T_S$ 线性且固定 $\mathbf1$，对词长归纳即可。
\end{proof}


\section{共同余数不能被平均改变}
从整数输入出发，分母只会增加 $p$ 的幂。因此当 $\ell\ne p$ 时，
所有中间数都仍然有良好的模 $\ell$ 意义：把分母中的 $p$ 替换成其模逆元即可。
若所有数的余数均为 $c$，平均值的余数就是 $p^{-1}(pc)=c$。
关键是本原化保证一个共同余数不可能在某个公因子处恰好全为零。

\begin{lemma}[非平均素数处的共同余数障碍]\label{lem:obstruction}
非零本原整数零和配置满足 $G\mid n$。若素数 $\ell\ne p$ 整除 $G$，则配置不可归零。
\end{lemma}
\begin{proof}
模 $G$ 的全部坐标同余于某个 $a$。本原性给出 $\gcd(a,G)=1$，零和给出 $na\equiv0\pmod G$，所以 $G\mid n$。若 $\ell\mid G$，共同余数模 $\ell$ 非零。所有中间值属于 $\Z[1/p]$，该环可约化至 $\F_\ell$；每次平均保持这个非零共同余数，故不能归零。
\end{proof}

称一个整数零和配置是\emph{合法的}，如果它在每个素数 $\ell\mid n$、$\ell\ne p$ 处不全同余。这个定义在未取本原代表时比 $G=p^e$ 更强；从本原输入出发二者相同。称保持合法性的操作为\emph{安全操作}。记
\[
d_p(n)=\#\{\ell:\ell\mid n,\ \ell\ne p\text{ 为素数}\}.
\]


\section{读懂最低p进赋值的论证}
对非零有理数 $x$，写 $x=p^h u/v$，其中 $p\nmid uv$，定义 $v_p(x)=h$，并令 $v_p(0)=+\infty$。
例如 $v_3(18)=2$，$v_3(2/9)=-2$。这里不需要完成化或 $p$ 进数域。
只用以下三条初等性质：
\[
v_p(xy)=v_p(x)+v_p(y),\qquad
v_p(x+y)\ge\min(v_p(x),v_p(y)),
\]
以及当两项的赋值不同时，后一不等式取等号。
最后一条可由提取较小的 $p$ 幂直接看出：一个模 $p$ 单位加上一个被 $p$ 整除的数仍是单位。

本题中特别有用的是：若 $1\le k<p$，则 $v_p(ka)=v_p(a)$。
所以一个最低值块被\emph{部分}选中时，不能靠别的更高赋值项抵消它。
所选和再除以 $p$，最低赋值还会下降一。这正是下一个引理的不变量。

\begin{lemma}[小维数中首次分数化不可挽回]\label{lem:fraction}
设 $p$ 为奇素数，$p\le n\le2p$。整数零和配置的一条成功平均路径不可能首次产生非整数。
\end{lemma}
\begin{proof}
首次非整数平均后，最低 $p$ 进赋值 $h<0$ 恰由一个等值的 $p$ 块承担。若下一选择与此块交于 $k$ 个位置，$1\le k<p$，则其和的赋值为 $h$，新均值的赋值为 $h-1$，并形成新的最低 $p$ 块。选整个最低块时操作恒等。若选择与之不交，则只能 $n=2p$ 且选中补集；零和迫使结果成为 $(a^p,(-a)^p)$。按引理\ref{lem:affine}，它等价于 $(1^p,(-1)^p)$，被模 $2$ 的共同非零余数阻碍。故最低负赋值既不能消失，也不能通过不交操作进入可解配置。
\end{proof}


\section{构造阈值下方的不可解输入}
现在选一个既含 $0$ 又含 $1$ 的配置，使 $G=1$ 自动成立，
同时令任何整数平均都几乎没有选择。负值 $-p$ 与 $p$ 个 $1$ 恰好平衡总和。

\begin{proposition}[位于2p的反例]\label{prop:lower}
对奇素数 $p$，配置 $(-p,0^{p-1},1^p)$ 本原、零和、$G=1$，但不可归零。
\end{proposition}
\begin{proof}
由引理\ref{lem:fraction}，成功路径必须始终采用整数平均。若一次选择包含 $j$ 份 $1$ 和 $e$ 份 $-p$，整性要求 $j-ep\equiv0\pmod p$，故 $j=0$ 或 $p$。唯一非恒等可能是选择 $(-p,0^{p-1})$，得到 $((-1)^p,1^p)$；引理\ref{lem:obstruction}排除该配置。
\end{proof}


\section{二平均的独立结论}
二平均同样属于素数平均，但它的最终阈值为 $4$。
这里直接引用完全混合的已发表刻画，并说明与本题符号的对应。
奇素数的证明不依赖这一引用。

二平均使用\cite{mixing}引言的完全混合刻画定理。其条件 MC 是：对每个正奇数 $b$，若整数配置模 $b$ 全同余，则均值也属于这个余数类。代入本原零和 $X$ 后，MC 等价于 $G$ 无奇素因子，故给出 $n\ge4$ 的判据。混合图按拓扑顺序执行时，每个二入二出节点可用两个原位置实现。另一方面，$(-3,1,2)$ 的 $G=1$，任意首次二平均产生 $(a,a,-2a)$、$a\ne0$；后继非恒等操作只把参数变成 $-a/2$，故三元不可归零。这证明推论\ref{cor:threshold}的二平均部分。

